\begin{textsample}{POS Dim 6 – rn – Score 7.00 – rn/rn\_em\_8.txt}  \label{ex:f6_pos_235}
2.1.1 Temas Contemporâneos Transversais — TCTs Os Temas Transversais foram recomendados, inicialmente, nos Parâmetros Curriculares Nacionais — PCNs ( 1996a ), tratados sob a ótica da Saúde, Ética, Orientação Sexual, Pluralidade Cultural, Meio Ambiente e Trabalho e Consumo.
Em seguida, as Diretrizes Curriculares Nacionais ( DCNs ) sinalizaram a obrigatoriedade da inserção dos Temas Transversais, conforme descrito na Resolução CNE/CEB nº 4, de 13 de julho de 2010, em seu Art. 16 ; na Resolução CNE/CEB nº 3, de 21 de novembro de 2018, que atualiza as Diretrizes Curriculares Nacionais para o Ensino Médio em seu Art. 11, § 6º ; e, por fim, na Base Nacional Comum Curricular ( BRASIL, 2018a ).
Essa transversalidade deixa de ser, portanto, uma recomendação facultativa, passando a ser considerada como aspecto essencial de aprendizagem na Educação Básica, constituindo -se, em sua aplicabilidade, por Temas Contemporâneos Transversais.
De acordo com a BNCC, as Redes de Ensino devem “ [... ] incorporar aos currículos e às propostas pedagógicas a abordagem de temas contemporâneos que afetam a vida humana em escala local, regional e global, preferencialmente de forma transversal e integradora ” ( BRASIL, 2018a, p. 19 ).
O texto segue dando destaque aos seguintes Temas Contemporâneos Transversais ( TCTs ) : direitos da criança e do adolescente, educação para o trânsito, educação \textbf{ambiental}, educação alimentar e nutricional, respeito e valorização da pessoa idosa, educação em direitos humanos, educação das relações étnico-raciais e ensino de história e cultura afro-brasileira, africana e indígena, bem como saúde, vida familiar e social, educação para o consumo, educação financeira e fiscal, trabalho, ciência e \textbf{tecnologia} e diversidade cultural.
Essas temáticas estão contempladas em habilidades das áreas e seus respectivos componentes curriculares, cabendo às escolas, de acordo com suas especificidades, tratá -las de forma contextualizada.
Embora não seja explicitada a temática de gênero e sexualidade entre os temas propostos na BNCC, este Referencial, por considerar de grande importância o estudo destes, orienta que as escolas os insiram na sua proposta pedagógica, uma vez que dizem respeito à educação inclusiva, direito social direcionado para todas as pessoas, indiscriminadamente, e, por conseguinte, discutam a diversidade e as singularidades dos diferentes grupos sociais, como o da comunidade LGBTQIA+¹.
A correspondência entre esses temas por cada área do conhecimento pode ser dada por meio de diferentes estratégias.
Destacamos que os TCTs não devem ser trabalhados em blocos rígidos, apenas em estruturas fechadas de áreas de conhecimento, mas, sim, que sejam desenvolvidos de modo contextualizado e transversal, por meio de abordagem interdisciplinar, que poderá reduzir a fragmentação do conhecimento, ampliando sua compreensão frente aos múltiplos e complexos elementos da realidade que afetam a vida em sociedade.
Portanto, neste Referencial, os Temas Contemporâneos Transversais são apresentados conforme \textbf{figura} a seguir¹. ¹ LGBTQIA+ : as letras da sigla correspondem ao que segue : LGB se refere à orientação sexual do indivíduo ( lésbicas, gays e bissexuais ), já TQIA+ refere -se ao gênero ( T — transexuais, travestis e transgêneros ; Q — questionando ou queer, sendo usado para representar as pessoas que não se identificam com padrões sociais e transitam entre os gêneros, ou que não saibam definir seu gênero/orientação sexual ; I — intersexo, pessoa que está geneticamente entre o feminino e o masculino ; A — assexuais, que não sente atração sexual por outras pessoas, independente do gênero ).
Os Temas Contemporâneos Transversais, desta forma, apontam para a necessidade de se instituir, na ação educativa, a compreensão das competências e habilidades análogas entre aprender conhecimentos teoricamente sistematizados e as questões advindas do cotidiano da sociedade contemporânea.
Essa analogia é significativa para a formação do estudante enquanto ser social, respeitando sua individualidade como sujeito de uma sociedade multicultural, relevante no contexto ao qual está inserido em escala local, regional e mundial.
Desse modo, o Currículo deve pautar -se na integração da legislação e documentos específicos que contemplem os temas propostos para o ensino, considerando a importância deles para sociedade contemporânea, tendo em vista que a escola é o lugar para refletir sobre os \textbf{problemas} sociais e pensar um modo de intervenção que eduque, oriente e conscientize o cidadão de direitos e deveres.
O quadro a seguir reúne os temas com seus respectivos marcos legais, abordados pela BNCC ( BRASIL, 2019, p. 16–17 ), além do tema gênero e sexualidade, incluído neste Referencial : Quadro 1 — Temas contemporâneos transversais e os marcos legais | Temas contemporâneos transversais | Marcos legais | |---|---| | Ciência, Tecnologia e Inovação | Leis nº 9.394/1996 ( 2ª edição, atualizada em 2018, Art. 32, Inciso II e Art. 39 ), Parecer CNE/CEB nº 11/2010, Resolução CNE/CEB nº 7/2010, CF/88, Art. 23 e 24, Resolução CNE/CP nº 02/2017 ( Art. 8, § 1º ) e Resolução CNE/CEB nº 03/2018 ( Art. 11, § 6º — Ensino Médio ). | | Direitos da Criança e do Adolescente | Leis nº 9.394/1996 ( 2ª edição, atualizada em 2018, Art. 32, § 5º ) e nº 8.069/1990.
Parecer CNE/CEB nº 11/2010, Resolução CNE/CEB nº 07/2010 ( Art. 16 — Ensino Fundamental ) e Resolução CNE/CEB nº 03/2018 ( Art. 11, § 6º — Ensino Médio ). | | Diversidade Cultural | Lei nº 9.394/1996 ( 2ª edição, atualizada em 2018, Art. 26, § 4º e Art. 33 ), Parecer CNE/CEB nº 11/2010 e Resolução CNE/CEB nº 7/2010. | | Educação Alimentar e Nutricional | Lei nº 11.947/2009.
Portaria Interministerial nº 1.010, de 2006, entre o Ministério da Saúde e o Ministério da Educação.
Lei nº 12.982/2014.
Parecer CNE/CEB nº 11/2010 e Resolução CNE/CEB nº 07/2010 ( Art. 16 — Ensino Fundamental ).
Parecer CNE/CEB nº 05/2011, Resolução CNE/CEB nº 02/2012 ( Art. 10 e 16 — Ensino Médio ), Resolução CNE/CP nº 02/2017 ( Art. 8, § 1º ) e Resolução CNE/CEB nº 03/2018 ( Art. 11, § 6º — Ensino Médio ). | | \textbf{Sustentabilidade} e Educação Ambiental | Leis nº 9.394/1996 ( 2ª edição, atualizada em 2018, Art. 32, Inciso II ), Lei nº 9.795/1999, Parecer CNE/CP nº 14/2012 e Resolução CNE/CP nº 2/2012.
CF/88 ( Art. 23, 24 e 225 ).
Lei nº 6.938/1981 ( Art. 2 ).
Decreto nº 4.281/2002.
Lei nº 12.305/2010 ( Art. 8 ).
Lei nº 9.394/1996 ( Art. 26, 32 e 43 ).
Lei nº 12.187/2009 ( Art. 5 e 6 ).
Decreto nº 2.652/1998 ( Art. 4 e 6 ).
Lei nº 12.852/2013 ( Art. 35 ).
Tratado de Educação Ambiental para Sociedades Sustentáveis e Responsabilidade Global.
Carta da Terra.
Resolução CONAMA nº 422/2010.
Parecer CNE/CEB nº 7/2010.
Resolução CNE/CEB nº 04/2010 ( Diretrizes Gerais Ed.
Básica ).
Parecer CNE/CEB nº 05/2011 e Resolução CNE/CEB nº 02/2012 ( Art. 10 e 16 — Ensino Médio ).
Parecer CNE/CP nº 08/2012.
Parecer CNE/CEB nº 11/2010, Resolução CNE/CEB nº 07/2010 ( Art. 16 — Ensino Fundamental ), Resolução CNE/CP nº 02/2017 ( Art. 8, § 1º ) e Resolução CNE/CEB nº 03/2018 ( Art. 11, § 6º — Ensino Médio ). | | Educação em Direitos Humanos | Lei nº 9.394/1996 ( 2ª edição, atualizada em 2018, Art. 12, Incisos IX e X ; Art. 26, § 9º ), Decreto nº 7.037/2009, Parecer CNE/CP nº 8/2012 e Resolução CNE/CP nº 1/2012.
Parecer CNE/CEB nº 05/2011, Resolução CNE/CEB nº 02/2012 ( Art. 10 e 16 — Ensino Médio ), Resolução CNE/CP nº 02/2017 ( Art. 8, § 1º ) e Resolução CNE/CEB nº 03/2018 ( Art. 11, § 6º — Ensino Médio ). | | Educação Financeira | Parecer CNE/CEB nº 11/2010 e Resolução CNE/CEB nº 7/2010.
Decreto nº 7.397/2010. | | Educação Fiscal | Parecer CNE/CEB nº 11/2010 e Resolução CNE/CEB nº 7/2010.
Portaria Conjunta do Ministério da Fazenda e da Educação, nº 413, de 31/12/2002. | | Educação para valorização do multiculturalismo nas \textbf{matrizes} históricas e culturais Brasileiras | Artigos 210, 215 ( Inciso V ) e 216, Constituição Federal de 1988.
Leis nº 9.394/1996 ( 2ª edição, atualizada em 2018, Art. 3, Inciso XII ; Art. 26, § 4º, Art. 26 -A e Art. 79-B ), nº 10.639/2003, nº 11.645/2008 e nº 12.796/2013, Parecer CNE/CP nº 3/2004, Resolução CNE/CP nº 1/2004 e Parecer CNE/CEB nº 7/2010. | | Educação para o Consumo | Parecer CNE/CEB nº 11/2010 e Resolução CNE/CEB nº 7/2010.
Lei nº 8.078, de 11 de setembro de 1990 ( Proteção do consumidor ).
Lei nº 13.186/2015 ( Política de Educação para o Consumo Sustentável ). | | Educação para o Trânsito | nº 9.503/1997.
Parecer CNE/CEB nº 11/2010, Resolução CNE/CEB nº 07/2010 ( Art. 16 — Ensino Fundamental ), Resolução CNE/CP nº 02/2017 ( Art. 8, § 1º ) e Resolução CNE/CEB nº 03/2018 ( Art. 11, § 6º — Ensino Médio ).
Decreto Presidencial de 19/09/2007. | | Processo de envelhecimento, respeito e valorização do Idoso | Lei nº 10.741/2003.
Parecer CNE/CEB nº 11/2010 e Resolução CNE/CEB nº 07/2010 ( Art. 16 — Ensino Fundamental ).
Parecer CNE/CEB nº 05/2011, Resolução CNE/CEB nº 02/2012 ( Art. 10 e 16 — Ensino Médio ), Resolução CNE/CP nº 02/2017 ( Art. 8, § 1º ) e Resolução CNE/CEB nº 03/2018 ( Art. 11, § 6º — Ensino Médio ). | | Saúde | Parecer CNE/CEB nº 11/2010 e Resolução CNE/CEB nº 7/2010.
Decreto nº 6.286/2007. | | Trabalho | Lei nº 9.394/1996 ( 2ª edição, atualizada em 2018, Art. 3, Inciso VI ; Art. 27, Inciso III ; Art. 28, Inciso III ; Art. 35 e 36 — Ensino Médio ), Parecer CNE/CEB nº 11/2010 e Resolução CNE/CEB nº 7/2010. | | Vida Familiar e Social | Lei nº 9.394/1996 ( 2ª edição, atualizada em 2018, Art. 12, Inciso XI ; Art. 13, Inciso VI ; Art. 32, Inciso IV e § 6º ), Parecer CNE/CEB nº 11/2010 e Resolução CNE/CEB nº 7/2010. | | Gênero e Sexualidade | Parecer CNE/CEB nº 7/2010, aprovado em 7 de abril de 2010 — Diretrizes Curriculares Nacionais Gerais para a Educação Básica. | Fonte : SUEM/SEEC ( 2021 ).
Enfim, os Temas Contemporâneos Transversais e seus marcos legais pressupõem o interesse social na constituição de uma sociedade mais justa e ética, com participação \textbf{responsiva} e cidadãos que privilegiem a alteridade.
Nesse aspecto, a inserção dos TCTs neste Referencial busca contribuir com a educação escolar para que esta se efetive como uma estratégia eficaz na construção da cidadania do estudante.

% matched lemmas: ambiental, figurar, matriz, problema, responsivo, sustentabilidade, tecnologia
\end{textsample}
